% GENERATED FILE. Edit canonical lessons, then run npm run generate:books.
% canonical-source-hash: fnv1a64:18e54d7547319d27
% canonical-lessons: UR-C191-chae-khana, UR-C191-girja-ghar, UR-C191-gurdvara, UR-C191-qabristan, UR-C191-minar

\chapter{Tea house, Church, Gurdwara, Graveyard, Minaret}
\label{ch:ur-tea-house-church-gurdwara-graveyard-minaret}
\hlchaptermodality{191}

\begin{chapteropening}
\textbf{By the end of this chapter:} \emph{I can say a tea house, a church, a gurdwara, a graveyard and a minaret.}

Complete the last lesson of chapter 191: I can say a tea house, a church, a gurdwara, a graveyard and a minaret.
\end{chapteropening}

\section[\texorpdfstring{\ur{چائے} \ur{خانہ} (chāe khānā)}{چائے خانہ (chāe khānā)}]{\ur{چائے} \ur{خانہ} (chāe khānā) — a tea house}

\label{lesson:UR-C191-chae-khana}



\begin{quote}
Before the new one: say the Urdu for a college, then the Urdu for a dwelling.
\end{quote}

\subsection*{You'll want to know: \ur{چائے} \ur{خانہ}}
\textbf{\ur{چائے} \ur{خانہ}} — \emph{chāe khānā} — \textquotedblleft{}a tea house\textquotedblright{}.

\begin{cousinweb}
\textbf{\ur{چائے} \ur{خانہ}} (\emph{chāe khānā}) is \textquotedblleft{}tea house\textquotedblright{}. \emph{chāe}, \textquotedblleft{}tea\textquotedblright{}, came through Persian \emph{chāy} from Chinese \emph{chá}, and \emph{khānā} is Persian \emph{khāna}, \textquotedblleft{}house, room\textquotedblright{}.
\end{cousinweb}

\begin{grammarlens}[title={What you've built}]
One more word for this chapter.
\end{grammarlens}

\subsection*{Your turn}
\begin{itemize}
\raggedright
  \item \emph{Say it:} \emph{chāe khānā}
  \item \emph{Say it:} \emph{chāe khānā}, once more
  \item \emph{Say it:} say \emph{chāe khānā}
  \item \emph{Recall:} say the Urdu for a college, then the Urdu for a dwelling, then say \emph{chāe khānā} again
\end{itemize}

\begin{tcolorbox}[breakable,colback=teal!4,colframe=teal!35!black,title={Before you move on}]
What does \ur{چائے} \ur{خانہ} mean? (\textquotedblleft{}A tea house\textquotedblright{}.) Say it once more.
\end{tcolorbox}
\section[\texorpdfstring{\ur{گرجا} \ur{گھر} (girjā ghar)}{گرجا گھر (girjā ghar)}]{\ur{گرجا} \ur{گھر} (girjā ghar) — a church}

\label{lesson:UR-C191-girja-ghar}



\begin{quote}
Before the new one: say the Urdu for a dwelling, then the Urdu for a tea house.
\end{quote}

\subsection*{You'll want to know: \ur{گرجا} \ur{گھر}}
\textbf{\ur{گرجا} \ur{گھر}} — \emph{girjā ghar} — \textquotedblleft{}a church\textquotedblright{}.

\begin{cousinweb}
\textbf{\ur{گرجا} \ur{گھر}} (\emph{girjā ghar}) is \textquotedblleft{}church house\textquotedblright{}. \emph{girjā} is usually traced to Portuguese \emph{igreja}, \textquotedblleft{}church\textquotedblright{}. \emph{ghar}, \textquotedblleft{}house\textquotedblright{}, is inherited; it is usually taken as a popular form of Sanskrit \emph{gṛha}, a link CDIAL calls phonetically difficult (CDIAL 4428).
\end{cousinweb}

\begin{grammarlens}[title={What you've built}]
One more word for this chapter.
\end{grammarlens}

\subsection*{Your turn}
\begin{itemize}
\raggedright
  \item \emph{Say it:} \emph{girjā ghar}
  \item \emph{Say it:} \emph{girjā ghar}, once more
  \item \emph{Say it:} say \emph{girjā ghar}
  \item \emph{Recall:} say the Urdu for a dwelling, then the Urdu for a tea house, then say \emph{girjā ghar} again
\end{itemize}

\begin{tcolorbox}[breakable,colback=teal!4,colframe=teal!35!black,title={Before you move on}]
What does \ur{گرجا} \ur{گھر} mean? (\textquotedblleft{}A church\textquotedblright{}.) Say it once more.
\end{tcolorbox}
\section[\texorpdfstring{\ur{گوردوارہ} (gurdvārā)}{گوردوارہ (gurdvārā)}]{\ur{گوردوارہ} (gurdvārā) — a gurdwara}

\label{lesson:UR-C191-gurdvara}



\begin{quote}
Before the new one: say the Urdu for a tea house, then the Urdu for a church.
\end{quote}

\subsection*{You'll want to know: \ur{گوردوارہ}}
\textbf{\ur{گوردوارہ}} — \emph{gurdvārā} — \textquotedblleft{}a gurdwara\textquotedblright{}.

\begin{cousinweb}
\textbf{\ur{گوردوارہ}} (\emph{gurdvārā}) is \textquotedblleft{}the guru's gateway\textquotedblright{}, the Sikh place of worship. It joins Sanskrit \emph{guru}, \textquotedblleft{}teacher\textquotedblright{}, and \emph{dvārā}, \textquotedblleft{}door, gateway\textquotedblright{}, from Sanskrit \emph{dvāra}, \textquotedblleft{}door\textquotedblright{}; CDIAL lists the inherited Hindi \emph{duār} under a form \emph{*duvāra} (CDIAL 6459).
\end{cousinweb}

\begin{grammarlens}[title={What you've built}]
One more word for this chapter.
\end{grammarlens}

\subsection*{Your turn}
\begin{itemize}
\raggedright
  \item \emph{Say it:} \emph{gurdvārā}
  \item \emph{Say it:} \emph{gurdvārā}, once more
  \item \emph{Say it:} say \emph{gurdvārā}
  \item \emph{Recall:} say the Urdu for a tea house, then the Urdu for a church, then say \emph{gurdvārā} again
\end{itemize}

\begin{tcolorbox}[breakable,colback=teal!4,colframe=teal!35!black,title={Before you move on}]
What does \ur{گوردوارہ} mean? (\textquotedblleft{}A gurdwara\textquotedblright{}.) Say it once more.
\end{tcolorbox}
\section[\texorpdfstring{\ur{قبرستان} (qabristān)}{قبرستان (qabristān)}]{\ur{قبرستان} (qabristān) — a graveyard}

\label{lesson:UR-C191-qabristan}



\begin{quote}
Before the new one: say the Urdu for a church, then the Urdu for a gurdwara.
\end{quote}

\subsection*{You'll want to know: \ur{قبرستان}}
\textbf{\ur{قبرستان}} — \emph{qabristān} — \textquotedblleft{}a graveyard\textquotedblright{}.

\begin{cousinweb}
\textbf{\ur{قبرستان}} (\emph{qabristān}) is \textquotedblleft{}place of graves\textquotedblright{}. It joins Arabic \emph{qabr}, \textquotedblleft{}a grave\textquotedblright{} (Lane, root \emph{q-b-r}), and the Persian ending \emph{-istān}, \textquotedblleft{}place\textquotedblright{}.
\end{cousinweb}

\begin{grammarlens}[title={What you've built}]
One more word for this chapter.
\end{grammarlens}

\subsection*{Your turn}
\begin{itemize}
\raggedright
  \item \emph{Say it:} \emph{qabristān}
  \item \emph{Say it:} \emph{qabristān}, once more
  \item \emph{Say it:} say \emph{qabristān}
  \item \emph{Recall:} say the Urdu for a church, then the Urdu for a gurdwara, then say \emph{qabristān} again
\end{itemize}

\begin{tcolorbox}[breakable,colback=teal!4,colframe=teal!35!black,title={Before you move on}]
What does \ur{قبرستان} mean? (\textquotedblleft{}A graveyard\textquotedblright{}.) Say it once more.
\end{tcolorbox}
\section[\texorpdfstring{\ur{مینار} (mīnār)}{مینار (mīnār)}]{\ur{مینار} (mīnār) — a minaret, a tower}

\label{lesson:UR-C191-minar}



\begin{quote}
Before the new one: say the Urdu for a gurdwara, then the Urdu for a graveyard.
\end{quote}

\subsection*{You'll want to know: \ur{مینار}}
\textbf{\ur{مینار}} — \emph{mīnār} — \textquotedblleft{}a minaret, a tower\textquotedblright{}.

\begin{cousinweb}
\textbf{\ur{مینار}} (\emph{mīnār}) is a loan from Arabic \emph{manār}, \textquotedblleft{}a lighthouse, a minaret\textquotedblright{}, which reached Urdu through Persian. Its Arabic root is \emph{n-w-r} (Lane). It is built on the same root as \emph{nūr}, \textquotedblleft{}light\textquotedblright{}.
\end{cousinweb}

\begin{grammarlens}[title={What you've built}]
One more word for this chapter.
\end{grammarlens}

\subsection*{Your turn}
\begin{itemize}
\raggedright
  \item \emph{Say it:} \emph{mīnār}
  \item \emph{Say it:} \emph{mīnār}, once more
  \item \emph{Say it:} say \emph{mīnār}
  \item \emph{Recall:} say the Urdu for a gurdwara, then the Urdu for a graveyard, then say \emph{mīnār} again
\end{itemize}

\begin{tcolorbox}[breakable,colback=teal!4,colframe=teal!35!black,title={Before you move on}]
What does \ur{مینار} mean? (\textquotedblleft{}A minaret, a tower\textquotedblright{}.) Say it once more.
\end{tcolorbox}
